\documentclass[11pt,a4paper]{article}
\usepackage[margin=2.5cm]{geometry}
\usepackage[T1]{fontenc}
\usepackage{multicol}   % two- or three-column lists
\usepackage{makeidx}
\makeindex

% ---------------------------------------------------------------
% MODULE I - Getting Started with LaTeX
% Syllabus units 1-4 (Text 1: Indian TeX Users Group primer,
% Tutorials I-V). Compile with: pdflatex, then bibtex/makeindex
% where a section below says so.
% ---------------------------------------------------------------

\title{Module I: Getting Started with \LaTeX}
\author{Study Notes -- MAT5FS112}
\date{}

\begin{document}
\maketitle
\tableofcontents
\newpage

% ===============================================================
\section{What is LaTeX?}
% ===============================================================

\LaTeX{} is a \emph{markup language} built on top of the \TeX{}
typesetting engine. You write a plain text \texttt{.tex} file that
describes the \emph{structure} of a document (title, sections,
lists, equations); the engine decides how it looks on the page.

The workflow is:
\begin{enumerate}
  \item Write source code in a \texttt{.tex} file (any text editor, or Overleaf).
  \item Compile it: \verb|pdflatex file.tex|.
  \item Open the resulting \texttt{file.pdf}.
\end{enumerate}

Compilation also creates helper files (\texttt{.aux}, \texttt{.log},
\texttt{.toc}). Always read the \texttt{.log} file when an error
appears; the first error is the one that matters.

% ===============================================================
\section{Anatomy of a document}
% ===============================================================

Every LaTeX file has two parts:
\begin{description}
  \item[Preamble] everything before \verb|\begin{document}|. It sets the
    document class, loads packages and defines global settings.
  \item[Body] everything between \verb|\begin{document}| and
    \verb|\end{document}|. This is the content you read.
\end{description}

\subsection{The smallest document}

\begin{verbatim}
\documentclass{article}
\begin{document}
Hello, world!
\end{document}
\end{verbatim}

\subsection{Document classes and options}

The class decides the overall look: \texttt{article} (papers,
assignments), \texttt{report} (long documents with chapters),
\texttt{book}, and \texttt{beamer} (slides, Module V). Options go in
square brackets:
\begin{verbatim}
\documentclass[11pt,a4paper]{article}
\end{verbatim}

\subsection{Packages}

Packages add features. Load them in the preamble:
\begin{verbatim}
\usepackage[margin=2.5cm]{geometry}   % page margins
\end{verbatim}

% ===============================================================
\section{Commands, environments and special characters}
% ===============================================================

A \emph{command} starts with a backslash: \verb|\textbf{bold}|.
An \emph{environment} is a block delimited by \verb|\begin{name}| and
\verb|\end{name}|. Environments must be closed in the reverse order in
which they were opened.

Some characters have special meaning and must be escaped:

\begin{center}
\begin{tabular}{ll}
\hline
You type & You get \\
\hline
\verb|\%| & \% \\
\verb|\$| & \$ \\
\verb|\&| & \& \\
\verb|\#| & \# \\
\verb|\_| & \_ \\
\verb|\{| and \verb|\}| & \{ and \} \\
\verb|\textbackslash| & \textbackslash \\
\verb|\textasciitilde| & \textasciitilde \\
\hline
\end{tabular}
\end{center}

A \texttt{\%} starts a comment: everything after it on that line is
ignored by the compiler.

% ===============================================================
\section{Document structure}
% ===============================================================

\subsection{Title block}

The title block is built from three commands in the preamble and
one in the body:
\begin{verbatim}
\title{My Title}
\author{Author Name}
\date{\today}          % or \date{} for no date
\maketitle             % in the body
\end{verbatim}

\subsection{Sectioning}

Sectioning commands number and title parts of the document:
\verb|\section|, \verb|\subsection|, \verb|\subsubsection|,
\verb|\paragraph|. Adding a star (\verb|\section*|) removes the number
and the table-of-contents entry.

\subsection{Table of contents}

\verb|\tableofcontents| prints the contents page. It needs
\emph{two} compilations: the first writes the entries to a
\texttt{.toc} file, the second reads them back.

\subsection{Lists}

\begin{multicols}{3}
\textbf{itemize} (bullets)
\begin{itemize}
  \item apples
  \item oranges
\end{itemize}
\end{multicols}

\begin{enumerate}
  \item First numbered item.
  \item Second numbered item.
    \begin{enumerate}
      \item Nested items are numbered differently.
    \end{enumerate}
\end{enumerate}

\begin{description}
  \item[Term] an explanation of the term.
\end{description}

% ===============================================================
\section{Bibliography}
% ===============================================================

\subsection{Hand-written bibliography (compiles as it is)}

For a short list, \verb|thebibliography| is enough. Each entry is
introduced by \verb|\bibitem| and cited with \verb|\cite|. Here the
citation \cite{lamport} and \cite{knuth} is resolved in this file.

\begin{thebibliography}{9}
  \bibitem{lamport} Leslie Lamport, \emph{\LaTeX: A Document
    Preparation System}, 2nd ed., Addison-Wesley, 1994.
  \bibitem{knuth} Donald E. Knuth, \emph{The \TeX book},
    Addison-Wesley, 1984.
\end{thebibliography}

\subsection{BibTeX database (for real projects)}

For larger projects, keep references in a separate \texttt{.bib}
file and let BibTeX format them. The database \texttt{refs.bib}
looks like this:
\begin{verbatim}
@book{lamport1994,
  author    = {Leslie Lamport},
  title     = {\LaTeX: A Document Preparation System},
  edition   = {2nd},
  publisher = {Addison-Wesley},
  year      = {1994}
}
\end{verbatim}

The document then uses:
\begin{verbatim}
\bibliographystyle{plain}
\bibliography{refs}        % refers to refs.bib, placed at the end
\end{verbatim}

Compile in this order: \verb|pdflatex|, \verb|bibtex|, \verb|pdflatex|,
\verb|pdflatex|.

% ===============================================================
\section{Index}
% ===============================================================

Mark index entries in the text with \verb|\index{term}|. Print the
index at the end with \verb|\printindex| (needs \verb|\makeindex| in
the preamble). Here we mark two terms:
\index{LaTeX}\index{preamble}

The index is produced in three steps: \verb|pdflatex|,
\verb|makeindex|, \verb|pdflatex|.

\printindex

\end{document}
